% Návrh sekce "Implementace oceňovacích metod" – binomické stromy
% Draft k projití a případným úpravám

\section{Implementace oceňovacích metod}

\subsection{Binomické stromy}

Pro implementaci binomických stromů byl zvolen Coxův-Rossův-Rubinsteinův model z balíčku \texttt{fOptions} \citep{fOptions}.

\subsubsection{Důvod použití hotové knihovny}

Vlastní implementace binomických stromů by vyžadovala nejen konstrukci samotného stromu, ale také ošetření okrajových případů (velmi krátké časy do expirace, extrémní strike ceny, speciální případy dividend). Knihovna \texttt{fOptions} poskytuje odladěnou implementaci CRR modelu jak pro evropské, tak pro americké opce, včetně korektního ošetření těchto okrajových situací. Z těchto důvodů bylo rozhodnuto použít tuto knihovnu, která byla stažena z GitHubu.

\subsubsection{Nastavení modelu}

Pro každou opci byl binomický strom sestaven s následujícími parametry:
\begin{itemize}
  \item \textbf{Počet kroků:} \(n = 100\). Tato hodnota představuje kompromis mezi přesností a výpočetní náročností. Hull \citep{Hull2021} uvádí, že pro praktické účely je 100 kroků dostačujících – cena opce se s dalším zvyšováním počtu kroků mění jen zanedbatelně.
  \item \textbf{Cena podkladového aktiva:} \(S\) – aktuální cena akcie v den sběru dat.
  \item \textbf{Realizační cena:} \(K\) – strike cena opce.
  \item \textbf{Bezriziková úroková míra:} \(r\) – interpolovaná hodnota z výnosové křivky amerických státních dluhopisů.
  \item \textbf{Volatilita:} \(\sigma\) – historická volatilita spočtená z denních uzavíracích cen.
  \item \textbf{Čas do expirace:} \(T\) – poměrná část roku.
  \item \textbf{Náklady na držení:} \(b = r\). Zde bylo provedeno zjednodušení – náklady na držení byly nastaveny rovny bezrizikové úrokové míře, což odpovídá scénáři bez dividendového výnosu. U dividendových akcií by korektní hodnota měla být \(b = r - q\); vliv tohoto zjednodušení je diskutován v sekci \ref{sec:porovnani}.
\end{itemize}

\subsubsection{Typy opcí}

Balíček \texttt{fOptions} rozlišuje typ opce parametrem \texttt{TypeFlag}:
\begin{itemize}
  \item \texttt{"ce"} -- evropská call opce,
  \item \texttt{"pe"} -- evropská put opce,
  \item \texttt{"ca"} -- americká call opce,
  \item \texttt{"pa"} -- americká put opce.
\end{itemize}

Pro evropské opce je cena vypočtena jako diskontovaná očekávaná výplata v rizikově neutrálním světě. Pro americké opce je v každém uzlu stromu porovnána pokračovací hodnota s vnitřní hodnotou při okamžitém uplatnění, čímž je korektně zohledněna možnost předčasného uplatnění.

\subsubsection{Způsob výpočtu}

Výpočet probíhal v jazyce R. Pro každou variantu opcí (americké call, americké put, evropské call, evropské put) byl zavolán \texttt{CRRBinomialTreeOption} na každém řádku datasetu. Díky použití balíčku \texttt{furrr} s paralelizací přes všechna dostupná jádra procesoru probíhalo toto zpracování paralelně, což výrazně zkrátilo celkový výpočetní čas.

Výsledná cena opce je vrácena jako S4 objekt; extrahována je pomocí slotu \texttt{@price}. Dataset byl poté rozšířen o sloupec \texttt{user\_crr} s cenami vypočtenými binomickým stromem.

\subsection{Blackův-Scholesův-Mertonův model}

Na rozdíl od binomických stromů nebyla pro BS-M model použita externí knihovna – model byl implementován vlastnoručně v jazyce R. Důvodem je jednak jednoduchost samotného vzorce (cena evropské opce je vyjádřena explicitně, není třeba iterativního výpočtu), jednak možnost přímé kontroly nad implementací Barone-Adesi-Whaley aproximace pro americké opce.

\subsubsection{Evropské opce}

Pro evropské call opce byl implementován vzorec Black-Scholes-Merton s dividendovým výnosem \(q\) \citep[s.~15.1]{McDonald2013}:

\[
\begin{aligned}
d_1 &= \frac{\ln(S/K) + (r - q + \frac{1}{2}\sigma^2)T}{\sigma\sqrt{T}}, \\[2pt]
d_2 &= d_1 - \sigma\sqrt{T}, \\[2pt]
C_{\text{eu}} &= S e^{-qT} N(d_1) - K e^{-rT} N(d_2),
\end{aligned}
\]

a pro evropské put opce:

\[
P_{\text{eu}} = K e^{-rT} N(-d_2) - S e^{-qT} N(-d_1).
\]

Implementace je přímá – funkce \texttt{bs\_european\_call} a \texttt{bs\_european\_put} přijímají parametry \(S, K, r, \sigma, T, q\) a vracejí cenu opce. Pro případy \(T \le 0\) (opce s expirací v den sběru dat) je cena určena jako vnitřní hodnota \(\max(S - K, 0)\) resp. \(\max(K - S, 0)\).

\subsubsection{Americké opce – aproximace Barone-Adesi-Whaley}

Barone-Adesi-Whaley aproximace rozšiřuje BS-M model o prémií za možnost předčasného uplatnění. Tato prémie je odvozena z kvadratické aproximace rozdílu mezi cenou americké a evropské opce \citep[Kapitola~3.1]{Haug2007}.

Klíčovým krokem je nalezení kritické ceny podkladového aktiva \(S^*\) (pro call) resp. \(S^{**}\) (pro put), která představuje hranici, při jejímž překročení je předčasné uplatnění optimální.

\paragraph{Kritická cena pro call opce.} Rovnice pro kritickou cenu \(S^*\) je:

\[
S^* - K = C_{\text{eu}}(S^*, K, \sigma, r, T, q) + \left(1 - e^{-qT} N[d_1(S^*)]\right) \frac{S^*}{q_2},
\]

kde \(q_2\) je kořen kvadratické rovnice odvozené z Black-Scholesovy parciální diferenciální rovnice. Tato rovnice je řešena Newton-Raphsonovou metodou (funkce \texttt{s\_star\_finder}). Jako počáteční odhad slouží kritická cena v nekonečném čase:

\[
S^*(\infty) = \frac{K}{1 - 2\left[-(N-1) + \sqrt{(N-1)^2 + 4M}\right]^{-1}},
\]

kde \(M = 2r / \sigma^2\) a \(N = 2(r - q) / \sigma^2\).

Samotná cena americké call opce je poté:

\[
C_{\text{am}} =
\begin{cases}
C_{\text{eu}}(S, K, \sigma, r, T, q) + A_2 \left(\frac{S}{S^*}\right)^{q_2} & \text{pro } S < S^*, \\[2pt]
S - K & \text{pro } S \ge S^*,
\end{cases}
\]

kde \(A_2 = \frac{S^*}{q_2} \left(1 - e^{-qT} N[d_1(S^*)]\right)\).

Pro call opce na akcie bez dividendy (\(q \le 0\)) je předčasné uplatnění nevýhodné a cena americké opce se rovná ceně evropské.

\paragraph{Kritická cena pro put opce.} Symetricky, pro put opce je hledána kritická cena \(S^{**}\) splňující:

\[
K - S^{**} = P_{\text{eu}}(S^{**}, K, \sigma, r, T, q) - \left(1 - e^{-qT} N[-d_1(S^{**})]\right) \frac{S^{**}}{q_1},
\]

kde \(q_1\) je záporný kořen kvadratické rovnice. Počáteční odhad je:

\[
S^{**}(\infty) = \frac{K}{1 - 2\left[-(N-1) - \sqrt{(N-1)^2 + 4M}\right]^{-1}}.
\]

Cena americké put opce:

\[
P_{\text{am}} =
\begin{cases}
P_{\text{eu}}(S, K, \sigma, r, T, q) + A_1 \left(\frac{S}{S^{**}}\right)^{q_1} & \text{pro } S > S^{**}, \\[2pt]
K - S & \text{pro } S \le S^{**},
\end{cases}
\]

kde \(A_1 = -\frac{S^{**}}{q_1} \left(1 - e^{-qT} N[-d_1(S^{**})]\right)\).

\subsubsection{Implementační detaily}

Newton-Raphsonova iterace pro nalezení \(S^*\) a \(S^{**}\) končí při relativní toleranci \(\alpha = 10^{-9}\) vztažené k strike ceně \(K\), tedy:

\[
\frac{|S^* - K - \text{RHS}(S^*)|}{K} < 10^{-9}.
\]

Pro zvýšení robustnosti je maximální počet iterací omezen na 1~000 (u \(S^*\)) a 10~000 (u \(S^{**}\)); v praxi konverguje algoritmus obvykle v řádu jednotek až desítek iterací.

Funkce pro americké opce jsou implementovány jako skalární a následně vektorizovány pomocí \texttt{Vectorize()} (proměnné \texttt{v\_bs\_american\_call} a \texttt{v\_bs\_american\_put}). To umožňuje efektivní dávkové zpracování celého datasetu.

Výpočet probíhal obdobně jako u binomických stromů – paralelizovaně přes všechny řádky datasetu. Výsledné ceny byly uloženy do sloupce \texttt{user\_bs}.

\subsection{Monte Carlo}

Metoda Monte Carlo byla implementována vlastnoručně v jazyce R. Na rozdíl od Black-Scholes-Mertonova modelu, který poskytuje explicitní vzorec, a na rozdíl od binomických stromů, kde byla použita hotová knihovna, zde bylo nutné implementovat jak generování náhodných cest, tak logiku rozhodování o předčasném uplatnění.

\subsubsection{Parametry simulace}

Pro všechny simulace byl zvolen jednotný počet \(N = 25\,000\) nezávislých realizací. Tato hodnota představuje kompromis mezi přesností odhadu a výpočetní náročností – směrodatná odchylka odhadu klesá s \(\frac{1}{\sqrt{N}}\), takže 25~000 simulací poskytuje dostatečnou přesnost při přijatelném výpočetním čase.

Pro generování náhodných veličin bylo použito rozdělení \(\mathcal{N}(0,1)\). Pro snížení rozptylu odhadu byla aplikována metoda antitetických proměnných: z každé sady \(N/2\) náhodných čísel \(Z_1, \dots, Z_{N/2}\) byla vytvořena druhá sada \(-Z_1, \dots, -Z_{N/2}\), čímž vzniká \(N\) párově záporně korelovaných realizací. Tím se rozptyl výsledného odhadu výrazně snižuje \citep[Kapitola~4.2]{Glasserman2004}.

\subsubsection{Evropské opce}

Pro evropské opce postačuje simulovat pouze koncovou cenu podkladového aktiva v čase expirace \(T\). Ta je generována geometrickým Brownovým pohybem:

\[
S_T = S \exp\!\left( (r - q - \tfrac{1}{2}\sigma^2)T + \sigma\sqrt{T}\,Z \right), \qquad Z \sim \mathcal{N}(0,1).
\]

Funkce \texttt{mc\_sim\_pathless} generuje vektor \(N\) koncových cen s využitím antitetických proměnných.

Cena evropské call opce je poté odhadnuta jako diskontovaný průměr výplat:

\[
C \approx e^{-rT} \frac{1}{N} \sum_{i=1}^{N} \max(S_T^{(i)} - K, 0),
\]

a obdobně pro put opci:

\[
P \approx e^{-rT} \frac{1}{N} \sum_{i=1}^{N} \max(K - S_T^{(i)}, 0).
\]

Tento způsob je výpočetně velmi efektivní, protože nevyžaduje ukládání celých cest.

\subsubsection{Americké opce – metoda Longstaff-Schwartz}

Pro americké opce je situace komplikovanější – v každém okamžiku je třeba rozhodnout, zda je výhodnější opci uplatnit, nebo pokračovat v držení. K tomuto účelu byla implementována metoda Longstaff-Schwartz \citep{LongstaffSchwartz2001}, která odhaduje pokračovací hodnotu pomocí regrese.

\paragraph{Generování cest.} Na rozdíl od evropských opcí je pro americké opce nutné simulovat celou trajektorii ceny. Čas do expirace \(T\) je rozdělen na týdenní kroky (\(\Delta t = 1/52\) roku). Počet kroků je tedy \(M = \max(1, \text{round}(T \cdot 52))\). Pro každý krok je cena generována postupnou aplikací GBM:

\[
S_{t+\Delta t} = S_t \exp\!\left( (r - q - \tfrac{1}{2}\sigma^2)\Delta t + \sigma\sqrt{\Delta t}\,Z \right).
\]

Funkce \texttt{mc\_sim\_pathed} vrací matici o rozměrech \(N \times M\), kde každý řádek představuje jednu simulovanou cestu. Antitetické proměnné jsou použity i zde.

\paragraph{Zpětná indukce s regresí.} Algoritmus postupuje od data expirace směrem k současnosti:

\begin{enumerate}
  \item V čase \(T\) je výplata opce \(h_T(S_T) = \max(S_T - K, 0)\) pro call, resp. \(\max(K - S_T, 0)\) pro put.
  \item Při přechodu z kroku \(k+1\) do kroku \(k\) jsou diskontované výplaty \(y = e^{-r\Delta t} \cdot \text{CF}_{k+1}\) regredovány na bázické funkce aktuální ceny \(S_k\). Jako bázické funkce byl zvolen kubický polynom \(\{1, X, X^2, X^3\}\), kde \(X = S/K\) je standardizovaná cena.
  \item Regrese je provedena pouze na cestách, kde je opce v penězích (ITM). Pro call opce je podmínka \(S_k > K\), pro put opce \(S_k < K\). Pokud je ITM cest méně než 100, regrese se v daném kroku přeskočí – kubická regrese na méně než 4~bodech by byla numericky nestabilní a i při 5--20~bodech by odhad koeficientů \(\beta\) měl příliš vysoký rozptyl, což by mohlo zkreslit rozhodování o předčasném uplatnění na všech cestách. V takovém případě jsou všechny výplaty diskontovány o \(e^{-r\Delta t}\) a algoritmus pokračuje předchozím krokem.
  \item Z regresních koeficientů \(\beta\) je vypočtena pokračovací hodnota \(\hat{C}(S_k) = \beta_0 + \beta_1 X + \beta_2 X^2 + \beta_3 X^3\) pro všechny cesty. Tato hodnota představuje odhad očekávané diskontované výplaty při optimálním budoucím rozhodování – tedy hodnotu opce, pokud v tomto kroku neuplatníme.
  \item Nyní porovnáme dvě hodnoty: okamžitou výplatu při uplatnění \(h_k(S_k)\) a pokračovací hodnotu \(\hat{C}(S_k)\). Pokud \(h_k(S_k) > \hat{C}(S_k)\), okamžitá výplata převyšuje očekávaný výnos z držení → opci uplatníme a CF nahradíme \(h_k(S_k)\). V opačném případě je výhodnější pokračovat v držení – CF ponecháme a diskontujeme o \(e^{-r\Delta t}\).
\end{enumerate}

Po dokončení zpětné indukce je cena opce odhadnuta jako průměr všech diskontovaných výplat napříč cestami:

\[
V_0 = \frac{1}{N} \sum_{i=1}^{N} \text{CF}_0^{(i)}.
\]

\subsubsection{Implementační detaily}

Stejně jako u předchozích metod probíhal výpočet paralelizovaně přes všechny řádky datasetu pomocí balíčku \texttt{furrr}. Každá opce byla zpracována samostatně na jednom jádře procesoru.

Pro zajištění reprodukovatelnosti byl nastaven seed generátoru náhodných čísel (\texttt{set.seed(42)}) před benchmarkovým vzorkováním. Samotné simulace využívají implicitní seedování balíčku \texttt{furrr} (parametr \texttt{.options = furrr\_options(seed = TRUE)}).

Výsledné ceny byly uloženy do sloupce \texttt{user\_mc}. Celkový výpočetní čas byl změřen pomocí \texttt{system.time} a uložen pro pozdější porovnání v sekci \ref{sec:porovnani}.

\subsection{Propojení výpočtu s daty}

\subsubsection{Struktura datasetu}

Dataset byl po načtení rozdělen do čtyř samostatných datových rámců podle stylu opce (americká/evropská) a typu (call/put). Toto rozdělení bylo provedeno pomocí filtrování:

\begin{verbatim}
euro_call_chains <- euro_chains[euro_chains$option_type == "Call", ]
euro_put_chains  <- euro_chains[euro_chains$option_type == "Put", ]
usa_call_chains  <- usa_chains[usa_chains$option_type == "Call", ]
usa_put_chains   <- usa_chains[usa_chains$option_type == "Put", ]
\end{verbatim}

Každý řádek datového rámce reprezentuje jednu opci se sloupci: \texttt{ticker}, \texttt{expiration}, \texttt{K} (strike), \texttt{S} (cena akcie), \texttt{r}, \texttt{sigma}, \texttt{T} (čas do expirace), \texttt{q} (dividendový výnos) a \texttt{option\_price} (tržní cena). Každé variantě metody přísluší jeden nový sloupec s vypočtenou cenou.

\subsubsection{Použití balíčku dplyr}

Všechny tři oceňovací metody byly aplikovány pomocí funkce \texttt{mutate()} z balíčku \texttt{dplyr} \citep{dplyr}, která přidává nové sloupce do datového rámce. Tento přístup zajišťuje, že vypočtené ceny zůstávají korektně přiřazeny ke svým opcím i při paralelním zpracování.

Pro Blackův-Scholesův-Mertonův model, který je vektorizovaný, stačilo volat funkci přímo uvnitř \texttt{mutate()}:

\begin{verbatim}
usa_call_chains <- usa_call_chains %>%
  mutate(user_bs = v_bs_american_call(S, K, r, sigma, T, q))
\end{verbatim}

Pro binomické stromy a Monte Carlo, kde každá opce vyžaduje samostatný výpočet, byla použita funkce \texttt{future\_map\_dbl()} z balíčku \texttt{furrr} \citep{furrr}. Ta aplikuje skalární funkci na každý řádek datového rámce paralelně napříč jádry procesoru:

\begin{verbatim}
usa_call_chains <- usa_call_chains %>%
  mutate(user_crr = future_map_dbl(
    row_number(),
    ~ CRRBinomialTreeOption(
        TypeFlag = "ca", S = S[.x], X = K[.x],
        Time = T[.x], r = r[.x], b = r[.x],
        sigma = sigma[.x], n = 100)@price
  ))
\end{verbatim}

Pro Monte Carlo byl navíc přidán parametr \texttt{.options = furrr\_options(seed = TRUE)} zajišťující reprodukovatelnost generátoru náhodných čísel.

\subsubsection{Měření výpočetního času}

Doba trvání každé metody byla změřena funkcí \texttt{system.time()}, která obaluje celé volání \texttt{mutate()} včetně paralelizace:

\begin{verbatim}
t_bs <- system.time({
  usa_call_chains <- usa_call_chains %>%
    mutate(user_bs = v_bs_american_call(S, K, r, sigma, T, q))
  # ... obdobně pro ostatní varianty
})
\end{verbatim}

Naměřené časy byly uloženy pro pozdější porovnání v sekci \ref{sec:porovnani}.

\subsubsection{Benchmark na 100 opcích}

Kromě dávkového zpracování celého datasetu byl proveden benchmark na náhodném vzorku 100 opcí z každé varianty. Vzorek byl fixován pomocí \texttt{set.seed(42)}. Každá metoda byla spuštěna samostatně na tomto vzorku a změřen její čas. Výsledky byly sestaveny do tabulky formátu \texttt{data.frame} s metodou, variantou opce, celkovým časem, průměrným časem na opci a počtem opcí za sekundu.

\begin{verbatim}
benchmark_100 <- data.frame(
  Variant = rep(c("American Call", "American Put", ...), each = 3),
  Method = rep(c("Black-Scholes", "CRR Binomial (n=100)",
                 "Monte Carlo (25k sims)"), times = 4),
  Total_Time_s = c(...),
  Avg_Time_Per_Option_s = c(...),
  Options_Per_Second = c(...)
)
\end{verbatim}

\subsubsection{Uložení výsledků}

Po dokončení výpočtů byly všechny čtyři datové rámce uloženy do samostatných CSV souborů ve složce \texttt{calculated\_prices/}:

\begin{verbatim}
write_csv(usa_call_chains, "./calculated_prices/american_calls.csv")
write_csv(usa_put_chains,  "./calculated_prices/american_puts.csv")
write_csv(euro_call_chains, "./calculated_prices/european_calls.csv")
write_csv(euro_put_chains,  "./calculated_prices/european_puts.csv")
\end{verbatim}

Tato data slouží jako vstup pro následující sekci věnovanou výsledkům ocenění.
